% Generated draft: requires mathematical and source review.
\documentclass[11pt]{article}
\usepackage[margin=25mm]{geometry}
\usepackage{fontspec}
\setmainfont{Latin Modern Roman}
\usepackage{amsmath,amssymb,mathtools,unicode-math}
\setmathfont{Latin Modern Math}
\usepackage{enumitem,needspace,xurl,hyperref}
\setlength{\parindent}{0pt}
\setlength{\parskip}{5pt}
\setlength{\emergencystretch}{3em}
\newcommand{\N}{\mathbb{N}}
\newcommand{\Z}{\mathbb{Z}}
\newcommand{\Q}{\mathbb{Q}}
\newcommand{\R}{\mathbb{R}}
\newcommand{\C}{\mathbb{C}}
\newcommand{\F}{\mathbb{F}}
\newcommand{\A}{\mathbb{A}}
\newcommand{\PP}{\mathbb{P}}
\newcommand{\E}{\mathbb{E}}
\newcommand{\eps}{\varepsilon}
\newcommand{\dd}{\,\mathrm{d}}
\begin{document}
\section*{UnsolvedMath ID: 7200064}
\textbf{Source problem code:} catalogue entry
\subsection{Definitions and mathematical statement}
A Euclidean $n$-simplex is the convex hull of $n+1$ affinely independent points in $\R^n$. An orthoscheme is an $n$-simplex whose vertices can be ordered $w_0,\ldots,w_n$ so that the consecutive edge vectors
\[
 w_1-w_0,\ w_2-w_1,\ldots,w_n-w_{n-1}
\]
are pairwise orthogonal.

Under the usual finite geometric-dissection convention, the question is whether
\[
 \forall n\ge1\ \forall S\text{ a Euclidean }n\text{-simplex},
 \quad\exists m\ge1\quad\exists Q_1,\ldots,Q_m
\]
with each $Q_i$ an $n$-dimensional orthoscheme contained in $S$, satisfying
\[
 S=\bigcup_{i=1}^m Q_i,\qquad
 \operatorname{int}(Q_i)\cap\operatorname{int}(Q_j)=\varnothing
 \quad(i\ne j).
\]
The pieces may have different shapes and sizes; the source gives no congruence requirement. This means actual nonoverlapping pieces, not a signed decomposition or an identity of formal volumes.

The short source wording does not specify finite versus infinite dissection, whether pieces must meet face to face, or any restriction on admissible cuts. The displayed finite-dissection interpretation is explicit for review; no extra face-to-face condition is imposed.
\subsection{Short English statement}
Can every Euclidean simplex, in every dimension, be cut into finitely many orthoschemes? The source summary reports affirmative low-dimensional and special-type cases, while treating the general case from dimension four as open.

Those reports mention Debrunner and an imprecise special facet-angle condition but give no exact theorem or reference. The summary says recent 2024--2026 progress was not fully verified. The pinned Wikipedia assessment remains subject to primary-source review, and the background's unpublished note conflicts with registry published metadata.
\subsection{Sources}
\begin{itemize}
\item Local UnsolvedMath registry, permanent ID 7200064, source code catalogue entry, ``Dissection into orthoschemes.'' Dataset curation and metadata: Ulam AI and UnsolvedMath Contributors, CC BY 4.0.
\item Wikipedia contributors, \emph{List of unsolved problems in mathematics}, Geometry, bullet 64, pinned revision 1366636767, \url{https://en.wikipedia.org/w/index.php?title=List_of_unsolved_problems_in_mathematics&oldid=1366636767#Geometry}; MediaWiki-source extraction, accessed 2026-07-29.
\item Debrunner's theorem and special-simplex results are named in the supplied summary without complete citations or exact hypotheses. No primary resolution reference is supplied.
\item Source-specific provenance: primary-source status audit NEEDS\_REVIEW and the background says unpublished pending that audit, while published metadata is true. Wikipedia text is CC BY-SA 4.0, with attribution and ShareAlike handling required for release, separately from dataset attribution.
\end{itemize}
\end{document}
